\chapter{Zákony zachování a pohyb v centrálním poli}

% \section{Integrály pohybu}
% Základní zákony zachování v mechanice bezprostředně souvisí se symetriemi časoprostoru (teorém Noetherové):
% \begin{itemize}
%     \item \textbf{Zachování mechanické energie}: plyne z homogenity času v konzervativním silovém poli $\vect{F} = -\nabla V(\vect{r})$:
%     \begin{equation}
%         E = T + V = \frac{1}{2}m\dot{\vect{r}}^2 + V(\vect{r}) = \text{konst.}
%     \end{equation}
%     \item \textbf{Zachování celkové hybnosti}: plyne z homogenity prostoru pro izolovanou soustavu částic.
%     \item \textbf{Zachování momentu hybnosti}: plyne z izotropie prostoru. Pro centrální pole $\vect{F} \parallel \vect{r}$ platí:
%     \begin{equation}
%         \vect{L} \coloneqq \vect{r} \crossprod \vect{p} = \text{konst.}
%     \end{equation}
% \end{itemize}

% \section{Pohyb v centrálním gravitačním poli}
% Vzhledem k zachování $\vect{L}$ probíhá pohyb v rovině kolmé na $\vect{L}$. V polárních souřadnicích $(r, \varphi)$ se úloha redukuje na jednorozměrný pohyb v efektivním potenciálu:
% \begin{equation}
%     E = \frac{1}{2}m\dot{r}^2 + V_{\text{eff}}(r), \qquad V_{\text{eff}}(r) = -\frac{\alpha}{r} + \frac{L^2}{2mr^2}.
% \end{equation}
